% cSpell:locale en,pt-BR - Adicionado para suporte ao corretor ortográfico
% Extensão Code Spell Checker no VS Code (muito útil!)
\chapter{Elementos pós-textuais}\label{cap:postextuais}

\section{Referências bibliográficas}\label{sec:bibliografia}

A lista de referências é impressa por \verb|\printbibliography|, chamado uma única vez, logo após o
último capítulo do texto. Sua elaboração segue a ABNT NBR 6023:2018~\cite{nbr6023:2018} --- que rege
a \emph{referência}, a descrição de cada obra --- e não a NBR 10520:2023~\cite{nbr10520:2023}, que
rege a \emph{citação}, a remissão feita no corpo do texto. As duas são tratadas na
\refcomp{}{sec:citacoes}.

\subsection{Crie o seu próprio arquivo de referências}\label{ssc:proprio-bib}

\textbf{O arquivo \texttt{referencias.bib} deste repositório é fictício.} Suas entradas descrevem
obras que não existem, e estão ali para servir aos exemplos de citação deste manual. Um trabalho
entregue apontando para ele tem bibliografia inventada --- e, ao contrário de uma lista vazia, esse
defeito não salta aos olhos de quem folheia o PDF: a lista parece perfeitamente legítima.

Crie o seu. É um arquivo de texto comum, com extensão \texttt{.bib}, na mesma pasta do arquivo
principal, e ele começa vazio. Declare-o no preâmbulo:

\begin{Verbatim}[fontsize=\small]
\usepackage[style=abnt]{biblatex}
\bibliography{minhas-referencias.bib}
\end{Verbatim}

Nos arquivos de partida as duas linhas acima vêm comentadas, e com elas a chamada da
bibliografia, exatamente para que ninguém aponte para o arquivo de exemplo sem perceber.
Descomente as três quando tiver o seu.

Gerenciadores como o Zotero e o Mendeley exportam diretamente para \texttt{.bib}. Se usar um deles,
revise o arquivo exportado: campos incompletos ou mal preenchidos produzem referências fora da
norma sem que nada falhe na compilação.

\section{Apêndices e anexos}\label{sec:apendices-anexos}

Os dois se distinguem pela \textbf{autoria}, e não pelo conteúdo. Apêndice é material elaborado pelo
próprio autor, para complementar sua argumentação. Anexo é documento de terceiro, incorporado como
fundamentação, prova ou ilustração.

Abra cada grupo uma vez e declare os elementos em seguida:

\begin{Verbatim}[fontsize=\small]
\apendices
\apendice{ape:questionario}{Questionário aplicado na pesquisa}
...

\anexos
\anexo{ane:certificado}{Certificado de participação no evento}
...
\end{Verbatim}

O comando \verb|\apendices| abre o grupo e produz uma entrada coletiva no sumário; cada
\verb|\apendice{<rótulo>}{<título>}| abre um apêndice, identificado por letra maiúscula consecutiva
--- A, B, C. Os comandos \verb|\anexos| e \verb|\anexo{}{}| funcionam do mesmo modo, com sequência
própria de letras. Ambos vêm depois do \verb|\backmatter|.

Ilustrações dentro de um apêndice ou anexo continuam na mesma sequência numérica do restante do
documento e aparecem normalmente nas listas. Este manual traz um exemplo de cada, ao final, com a
distinção de autoria explicada no próprio texto.

\section{Abreviaturas, siglas e símbolos}\label{sec:abreviaturas-siglas-simbolos}

Caso seu documento utilize muitas abreviaturas, siglas ou símbolos técnicos, é recomendável incluir
uma lista dedicada a esses elementos. Isso ajuda o leitor a entender rapidamente o significado de
termos técnicos ou abreviados usados ao longo do texto. Também serve como referência rápida, caso o
leitor não se lembre imediatamente do significado de uma sigla ou símbolo específico. A lista de
abreviaturas pode ser construída manualmente ou utilizando pacotes específicos do \LaTeX, como o
\texttt{nomencl} ou o \texttt{glossaries}. Este modelo inclui um arquivo de exemplo para a lista de
abreviaturas e siglas, localizado em \texttt{Preambulo/lista-abreviaturas.tex}. Para
incluir a lista no seu documento, descomente a linha correspondente no arquivo principal do
documento. Para imprimir a lista de abreviaturas e a lista de símbolos, descomente os comandos
\verb|\printglossary|, no arquivo principal do modelo. A lista de abreviaturas e a lista de símbolos
serão geradas automaticamente com base nas definições fornecidas no arquivo de abreviaturas e no
arquivo de símbolos.

Para referenciar uma abreviatura ou símbolo no texto, utilize o comando \verb|\gls{}|, passando como
argumento a chave definida para a abreviatura ou símbolo no arquivo de definições. Por exemplo, para
referenciar a abreviatura \emph{Mínimo Múltiplo Comum}, definida com a chave \texttt{mmc}, utilize o
comando \verb|\gls{mmc}|. Veja: \gls{mmc}. Na primeira vez que a abreviatura for referenciada, o
nome completo será exibido junto com a abreviatura entre parênteses. Nas referências subsequentes,
apenas a abreviatura será exibida, observe: \gls{mmc}. Isso garante que o leitor compreenda o
significado da abreviatura desde a primeira menção. A lista de abreviaturas e símbolos será
atualizada automaticamente com base nas referências feitas no texto, garantindo apenas abreviaturas
que foram citadas no texto apareçam na lista. Um outro ponto importante, é que o pacote
\texttt{glossaries} é integrado ao pacote \texttt{hyperref}, garantindo que todas as referências a
abreviaturas e símbolos sejam \emph{hyperlinks} clicáveis, facilitando a navegação no documento.

Símbolos funcionam de forma parecida, eles devem ser citados no texo para que apareçam na lista de
símbolos. Use o comando \verb|\gls{}| também para referenciar os símbolos, veja: \gls{symb:pi}. Caso
deseje imprimir a descrição do símbolo no texto, utilize o comando \verb|\glsdesc{}|, veja a
descrição para o símbolo \gls{symb:pi}: \glsdesc{symb:pi}. Se seu trabalho possui muitas equações
usando símbolos específicos, é uma boa prática mantê-los organizados em um arquivo de definições
separado, como o fornecido neste modelo em \texttt{Preambulo/lista-simbolos.tex}.

Um ponto importante: a lista de abreviaturas não é gerada automaticamente. Ao incluir uma nova
abreviatura ou símbolo no arquivo de definições, você deve compilar o documento para gerar o
documento auxiliar necessário para a lista de abreviaturas. Dependendo do ambiente de compilação que
você está utilizando, pode ser necessário executar comandos adicionais para gerar a lista de
abreviaturas corretamente. Consulte a documentação do pacote \texttt{glossaries}
(\url{https://www.ctan.org/pkg/glossaries}) para mais detalhes sobre como gerar a lista de
abreviaturas e a lista de símbolos. Geralmente, isto envolve uma chamada ao comando
\texttt{makeglossaries} passando o nome do arquivo principal na linha de comando.

\section{As listas, e quando não chamá-las}\label{sec:listas}

Quatro listas de ilustrações estão disponíveis, além das listas de abreviaturas e de símbolos:

\begin{itemize}
    \item \verb|\listoffigures| --- a lista de ilustrações da norma, impressa por este modelo
          sob o título ``LISTA DE FIGURAS'';
    \item \verb|\listoftables| --- lista de tabelas;
    \item \verb|\listofquadros| e \verb|\listofgraficos| --- listas próprias dos tipos Quadro e
          Gráfico. Um tipo criado por \verb|\novotipoilustracao| ganha a sua automaticamente.
\end{itemize}

Todas são elementos \textbf{opcionais}, e devem ser chamadas apenas quando houver o que listar.

\textbf{Uma lista vazia não deixa de ser emitida: ela sai como uma página de título seguida de
nada.} O modelo não verifica se há itens antes de imprimir o cabeçalho. É por isso que, nos
arquivos de partida, as quatro chamadas vêm comentadas, com a explicação ao lado --- um trabalho
recém-começado não tem ilustração alguma, e a primeira compilação produziria uma página com o
título ``LISTA DE FIGURAS'' e o resto em branco.

Descomente cada uma quando o documento passar a ter o elemento correspondente. O mesmo vale para os
glossários: \verb|\printglossary| sem nenhuma sigla citada no texto produz o mesmo tipo de página
órfã.

Os títulos das duas listas de glossário também são redefiníveis, no mesmo molde dos nomes tratados
na \refcomp{}{sec:nomes-personalizaveis}: \texttt{\textbackslash acronymlistname} para a lista de
abreviaturas e siglas, \texttt{\textbackslash symbolslistname} para a de símbolos. Ambos
acompanham o idioma do \texttt{babel}, e ambos são passados explicitamente na chamada, o que
permite trocar o título de uma chamada só sem redefinir o comando:

\begin{Verbatim}[fontsize=\small]
\printglossary[type=\acronymtype,title={\acronymlistname}]
\end{Verbatim}
